\begin{frame}{Literature lineage}
\framesubtitle{Why artefact-level LLM measurement matters}

\begin{itemize}
    \item Classic evaluation methods were limited to surface-level overlap \citep{papineni2001bleu,lin2004rouge} or semantic similarity \citep{zhang2019bertscore, zhao2019moverscore}
    % \item LLMs solved these problems by acting as \alert{general-purpose evaluators} \citep{liu_g-eval_2023, zheng_judging_2023}
    \item LLMs addressed these limitations by acting as \textbf{general-purpose evaluators} \citep{liu_g-eval_2023, zheng_judging_2023}
    \item Current frontier: automatically evaluating text on multiple dimensions \citep{hashemi_llm-rubric_2024,kolaSAJASimpleApproach2026,bang_transforming_2025}
    \begin{itemize}
        \item This adheres to the multi-dimensional nature of text quality \citep{fabbri_summeval_2021}
    \end{itemize}
\end{itemize}

\vspace{0.1cm}

\begin{block}{Core problem}
    Multi-dimensional LLM evaluation produces criterion-level scores, but it is often unclear how these scores should be recomposed into a meaningful claim about artefact quality.
\end{block}

% \vspace{0cm}

\begin{exampleblock}{Thesis response}
    Define artefact evaluation as a measurement problem: quality dimensions are theoretically specified, measured separately with LLMs, and recomposed to collect validity evidence for the resulting construct-level measurement claim.
\end{exampleblock}

\end{frame}